\documentclass[12pt,a4paper]{article}
\usepackage{amsmath, amssymb}
\usepackage{booktabs}
\usepackage{geometry}
\usepackage{parskip}
\geometry{margin=2.5cm}

\title{Exercise: Binomial Option Pricing}
\author{}
\date{}

\begin{document}
\maketitle

\section*{Setup}

We have a two-state, one-period model where the current stock price is $S = 100$. It can either go up to $S_u = 150$ or down to $S_d = 75$. The interest rate is $r = 0$ and the real-world probability of the up-move is $p = 1/2$. We look at a European call option with strike $K = 110$.

First, let's compute the option payoffs in each state:
\[
C_u = \max(150 - 110, 0) = 40, \qquad C_d = \max(75 - 110, 0) = 0.
\]

\section*{Part 1: Pricing with real-world probabilities vs.\ risk-neutral valuation}

\subsection*{Using real-world probabilities}

A first (naive) approach would be to just take the expected payoff under the real-world probability $p = 1/2$ and discount it. Since $r = 0$, discounting doesn't change anything here, so we get:
\[
C^{\text{real}} = p \cdot C_u + (1-p) \cdot C_d = \frac{1}{2} \cdot 40 + \frac{1}{2} \cdot 0 = 20.
\]

\subsection*{Using risk-neutral valuation}

The correct approach is to use the risk-neutral probability $q$, which is found by requiring that the stock price is a martingale (i.e., it grows at the risk-free rate). With $r = 0$ this means:
\[
S = q \cdot S_u + (1-q) \cdot S_d.
\]
Plugging in:
\[
100 = 150q + 75(1-q) = 75 + 75q \implies q = \frac{25}{75} = \frac{1}{3}.
\]

Now we price the call under $q$:
\[
C^{\text{RN}} = q \cdot C_u + (1-q) \cdot C_d = \frac{1}{3} \cdot 40 + \frac{2}{3} \cdot 0 \approx 13.33.
\]

\subsection*{Comparison}

The two methods give different prices: $C^{\text{real}} = 20$ vs.\ $C^{\text{RN}} = 13.33$. The reason is that the real-world probability $p$ has nothing to do with the no-arbitrage price of the option --- it reflects what investors think will happen, but that is already priced into the current stock price $S = 100$. The risk-neutral price of $13.33$ is the only price consistent with no-arbitrage.

\section*{Part 2: Arbitrage strategy}

Now suppose someone sells us the call at the real-world price of $20$. Since the true no-arbitrage value is only $13.33$, the option is overpriced and we can exploit this.

\subsection*{Finding the replicating portfolio}

We want to find a portfolio $(\Delta, B)$ --- $\Delta$ shares of the stock and $B$ in cash (negative meaning we borrow) --- that exactly replicates the call payoffs:
\[
\begin{cases}
150\Delta + B = 40 \\[4pt]
75\Delta + B = 0
\end{cases}
\]
Subtracting the second equation from the first gives $75\Delta = 40$, so:
\[
\Delta = \frac{40}{75} = \frac{8}{15} \approx 0.53.
\]
And from the second equation: $B = -75 \cdot \frac{8}{15} = -40$ (so we borrow 40).

The cost of setting up this replicating portfolio today is:
\[
\Delta \cdot S + B = \frac{8}{15} \cdot 100 - 40 = 53.33 - 40 = 13.33,
\]
which is exactly the risk-neutral price --- good, this is consistent.

\subsection*{The arbitrage}

Since the call is sold to us at $20$ but only costs $13.33$ to replicate, we can do the following:

\medskip
\begin{center}
\begin{tabular}{lccc}
\toprule
\textbf{Position} & $t = 0$ & $t = 1$ (up) & $t = 1$ (down) \\
\midrule
Buy replicating portfolio & $-13.33$ & $+40$ & $0$ \\
\quad (buy $\frac{8}{15}$ shares, borrow 40) & & & \\
Sell the call (at real-world price) & $+20.00$ & $-40$ & $0$ \\
\midrule
\textbf{Net} & $\mathbf{+6.67}$ & $\mathbf{0}$ & $\mathbf{0}$ \\
\bottomrule
\end{tabular}
\end{center}

\medskip
We pocket $\mathbf{6.67}$ at $t = 0$ and have no future obligations in either state. This is a pure arbitrage profit. The fact that such a strategy exists shows that pricing the option with real-world probabilities leads to an inconsistency with no-arbitrage, confirming that the risk-neutral price of $13.33$ is the correct one.

\end{document}